\documentclass{article}
\usepackage{amsmath}
\usepackage{parskip}
\usepackage{url}
\usepackage[brazilian]{babel}
\title{Clube de Matemática --- Reunião 02}
\author{Herom Silva}
\date{Maio de 2026}

\begin{document}

\maketitle

\section{Introdução}

Neste documento, apresentamos como a segunda reunião do Clube de Matemática 
é estruturado, junto com o problema central que guiará nossa discussão.
Primeiro definiremos algumas notações que usaremos para enunciar o problema, 
junto com possíveis técnicas para resolvê-lo, uma generalização do problema 
e alguns recursos para explorá-lo com maior profundidade.
\newpage
\section{Notação}

Aqui definiremos a notação que utilizaremos no problema. Em livros de Lógica 
padrão veremos $P$ como uma proposição, isto é, uma sentença declarativa com 
um valor de verdade definido como $Verdade$ ou $Falso$.

Dizemos que $P$ tem valor $Verdade$ quando a proposição definida por $P$ 
é verdadeira, sendo um exemplo
$$ P = \text{A Terra é um planeta} $$
$P$ é $Verdade$ porque a Terra é um planeta.

Dizemos que $P$ tem valor $Falso$ quando a proposição definida por $P$ 
é falsa, sendo um exemplo
$$ P = \text{A Lua é feita de queijo}$$
$P$ é $Falso$ porque a Lua não é feita de queijo.

Podemos negar o valor de uma proposição pelo símbolo $\neg$ para obter 
seu valor inverso. Como exemplo
$$ P = \text{A Terra é um planeta}$$
$$ \neg P = \text{A Terra não é um planeta}$$
Pode-se criar uma nova proposição usando operadores lógicos como $\land$ 
e $\lor$, que representam respectivamente o gramatical ``e'', ``ou''. 
Como exemplo, podemos criar as seguintes proposições
$$P = \text{A Terra é um planeta}$$
$$Q = \text{Ela tem água}$$
$$ P \land Q = \text{A Terra é um planeta e ela tem água} = Verdade$$
$$P \lor Q = \text{A Terra é um planeta ou ela tem água} = Verdade$$

Para uma proposição usando o operador $\land$ ser verdadeira, ambos $P$ 
e $Q$ devem ser verdadeiros, e para o $\lor$ precisamos que $P$ ou $Q$ 
sejam verdadeiros.

Outros operadores conhecidos como operadores conectivos são $\implies$ e 
$\iff$, eles representam respectivamente o gramatical ``se então'' e 
``se e somente se''. Como exemplo podemos criar as seguintes proposições
$$P = \text{Está chovendo}$$
$$Q = \text{O chão está molhado}$$
$$ P \implies Q = \text{Se está chovendo então o chão está molhado} = Verdade$$
$$P \iff Q = \text{Está chovendo se e somente se o chão está molhado} = Falso$$
Para uma proposição usando o operador $\implies$ ser verdadeira, $P$ deve 
ser $Falso$ ou $Q$ deve ser $Verdade$, e para o $\iff$ precisamos que 
$P$ e $Q$ tenham o mesmo valor de verdade.
\newpage
\section{Técnicas}

Uma das técnicas comuns usadas em lógica proposicional é a tabela-verdade, 
uma representação tabular dos valores que uma proposição pode assumir dado 
os valores de verdade de seus componentes. A tabela abaixo ilustra isso 
para as proposições $P \land Q$, $P \lor Q$, $P \implies Q$, $P \iff Q$ 
e $\neg P$

\begin{center}
\begin{tabular}{|c c|c|c|c|c|c|}
\hline
$P$ & $Q$ & $P \land Q$ & $P \lor Q$ & $P \implies Q$ & $P \iff Q$ & $\neg P$\\
\hline
V & V & V & V & V & V & F \\
V & F & F & V & F & F & F \\
F & V & F & V & V & F & V \\
F & F & F & F & V & V & V \\
\hline
\end{tabular}
\end{center}

Outra técnica comum é o uso de equivalências lógicas, isto é, o uso de 
regras conhecidas para simplificar ou transformar uma proposição em uma 
equivalente. Duas proposições são ditas logicamente equivalentes quando 
compartilham os mesmos valores de verdade para todas as combinações de 
$P$ e $Q$, denotado $P \equiv Q$. Algumas equivalências comuns são:
$$\neg(\neg P) \equiv P \quad \text{(Dupla Negação)}$$
$$\neg(P \land Q) \equiv \neg P \lor \neg Q \quad \text{(Lei de De Morgan)}$$
$$\neg(P \lor Q) \equiv \neg P \land \neg Q \quad \text{(Lei de De Morgan)}$$
$$P \implies Q \equiv \neg P \lor Q \quad \text{(Equivalência da Implicação)}$$
\newpage
\section{Problema}
Dada a proposição $\neg(P \land Q)$, conhecida como operador NAND e 
denotada $P \uparrow Q$, cuja tabela-verdade é

\begin{center}
\begin{tabular}{|c c|c|}
\hline
$P$ & $Q$ & $P \uparrow Q$\\
\hline
V & V & F \\
V & F & V \\
F & V & V \\
F & F & V \\
\hline
\end{tabular}
\end{center}

Você consegue expressar os operadores $\neg$, $\land$ e $\lor$ 
usando apenas $\uparrow$?
\newpage
\section{Generalização}
O operador NAND não é o único operador com propriedades interessantes. 
O operador NOR, denotado $P \downarrow Q$ e definido como $\neg(P \lor Q)$, 
é outro operador que pode ser construído a partir dos operadores que vimos. 
Sendo sua tabela-verdade

\begin{center}
\begin{tabular}{|c c|c|}
\hline
$P$ & $Q$ & $P \downarrow Q$\\
\hline
V & V & F \\
V & F & F \\
F & V & F \\
F & F & V \\
\hline
\end{tabular}
\end{center}

Um conjunto de operadores é dito \textit{funcionalmente completo} se toda 
proposição puder ser expressa usando apenas operadores desse conjunto. 
Por exemplo, o conjunto $\{\neg, \land\}$ é funcionalmente completo, 
pois podemos expressar $\lor$ usando a Lei de De Morgan:
$$P \lor Q \equiv \neg(\neg P \land \neg Q)$$

Um resultado ainda mais surpreendente é que um único operador pode ser 
suficiente para expressar toda a lógica proposicional. Tal operador é 
chamado de \textit{operador funcionalmente completo}. Uma pergunta natural 
a se fazer é se o operador NOR compartilha essa propriedade com o NAND, 
isto é, podemos expressar $\neg$, $\land$ e $\lor$ usando apenas $\downarrow$?
\newpage
\section{Recursos}

Os recursos a seguir são recomendados para explorar a lógica proposicional 
e a completude funcional em maior profundidade.

\begin{thebibliography}{9}

\subsection*{Nível Introdutório}

\bibitem{mortari2001}
Mortari, C. A.
\textit{Introdução à Lógica}.
UNESP, 2001.

\bibitem{copi2011}
Copi, I. M.; Cohen, C.; McMahon, K.
\textit{Introduction to Logic}, 14th ed.
Pearson, 2011.

\bibitem{priest2008}
Priest, G.
\textit{Logic: A Very Short Introduction}.
Oxford University Press, 2008.

\subsection*{Nível Intermediário}

\bibitem{alencar2002}
Alencar Filho, E.
\textit{Iniciação à Lógica Matemática}, 21ª ed.
Nobel, 2002.

\bibitem{enderton2001}
Enderton, H. B.
\textit{A Mathematical Introduction to Logic}, 2nd ed.
Academic Press, 2001.

\bibitem{barak2023}
Barak, B.
\textit{Introduction to Theoretical Computer Science}.
Disponível em: \url{https://introtcs.org/public/}

\subsection*{Nível Avançado}

\bibitem{mendelson2015}
Mendelson, E.
\textit{Introduction to Mathematical Logic}, 6th ed.
CRC Press, 2015.

\bibitem{machover1996}
Machover, M.
\textit{Set Theory, Logic and Their Limitations}.
Cambridge University Press, 1996.

\end{thebibliography}

\end{document}